\section{Transport interpretations and finite-support
transfer}\label{app:finite-support-transfer}

\subsection{Valuations and adjustment costs}
\label{sec:transport-economic-interpretation}

For the economic interpretations, consider probability laws $\mu,\nu$ with
compact support in Euclidean space and a continuous nonnegative transport
cost $c(x,y)$.
Write $\Pi(\mu,\nu)$ for their couplings.
Kantorovich duality equates the least allocation cost with the greatest
supporting valuation \citep{panaretos_invitation_2020}:
\begin{equation}\label{eq:transport-economic-dual}
  \inf_{\pi\in\Pi(\mu,\nu)}\int c(x,y)\,d\pi
  =\sup_{\varphi(x)+\psi(y)\le c(x,y)}
  \left\{\int\varphi\,d\mu+\int\psi\,d\nu\right\}.
\end{equation}
The supremum ranges over continuous origin and destination potentials on
the compact supports.
The inequality makes their combined valuation a lower bound on every
feasible shipment cost; optimizing these valuations attains the primal value.
For metric cost $c(x,y)=\norm{x-y}$, the Kantorovich--Rubinstein
specialization uses a single valuation function:
\begin{equation}\label{eq:transport-economic-w1}
  W_1(\mu,\nu)
  =\sup_{|f(x)-f(y)|\le\norm{x-y}}
  \left\{\int f\,d\mu-\int f\,d\nu\right\}.
\end{equation}
The constraint is the spatial-arbitrage bound for symmetric unit transport
costs.
Quadratic cost uses the two-potential dual in
\eqref{eq:transport-economic-dual}.
These potentials value transport constraints; $L$ in the main text bounds
the separate transmission of characteristics into exposures.

The dynamic interpretation of quadratic cost considers paths of probability
laws $\mu_t$ and velocity fields $v_t$ in the ambient Euclidean space.
The Benamou--Brenier formula is
\begin{equation}\label{eq:transport-economic-action}
  W_2^2(\mu,\nu)
  =\inf_{\mu_t,v_t}\int_0^1\int\norm{v_t(x)}^2\,d\mu_t(x)\,dt,
  \qquad
  \partial_t\mu_t+\nabla\!\cdot(v_t\mu_t)=0,
\end{equation}
with endpoints $\mu_0=\mu$ and $\mu_1=\nu$, and the continuity equation
understood weakly \citep{benamou_brenier_2000}.
An instantaneous cost $\norm{v}^2/2$ therefore gives minimum cost $W_2^2/2$.
Unit-normalized endpoints retain this interpretation with paths in ambient
space, consistent with the empirical chord metric.

To make the curvature connection precise, let an adjustment penalty $a(v)$
be twice continuously differentiable near its minimum at $v=0$, with
positive-definite Hessian $Q$.
Then $a(v)-a(0)=\tfrac12 v^\top Qv+o(\norm{v}^2)$.
The change of coordinates $z=Q^{1/2}x$ expresses the local quadratic
penalty in Euclidean units.
This supplies the adjustment-cost interpretation; the paper's exposure
geometry is fixed separately by the common-shock covariance operator $\Gamma$.

\subsection{Transfer for finite characteristic clouds}

The metric bracket in \Cref{cor:random-w2-metric} supplies the main-text
covariance restriction.
The following complementary bound compares squared transport costs directly
for indexed finite clouds with matching atom masses and a bounded support diameter.

\begin{proposition}[Finite-support transfer]\label{prop:random-w2-transfer}
  Write each finite characteristic cloud as
  \[
    \widehat C_i=\sum_a p_{ia}\delta_{x_{ia}},
    \qquad
    \widehat C_j=\sum_b p_{jb}\delta_{x_{jb}}.
  \]
  Fix one transmission draw $u_{ia}$ for each characteristic atom and define
  the indexed exposure clouds by
  \[
    \widehat P_i:=\sum_a p_{ia}\delta_{T_i(x_{ia},u_{ia})},
    \qquad
    \widehat P_j:=\sum_b p_{jb}\delta_{T_j(x_{jb},u_{jb})}.
  \]
  $G_{ij}:=W_{2,d_\Omega}(\widehat C_i,\widehat C_j)$ and
  $\tau_{ij}:=\tau_i+\tau_j$. Let $D_{ij}$ be any maintained bound on the
  risk-coordinate support diameter, that is on
  $\norm{t_\Gamma(x)-t_\Gamma(y)}$ and on $\norm{Z_i-Z_j}$ across all pairs of
  indexed support points. Then, under \Cref{ass:random-common-map},
  \begin{align}
    L^{-2}G_{ij}^{2}-2\tau_{ij}D_{ij}-\tau_{ij}^{2}
    &\le W_{2,\Gamma}^{2}(\widehat P_i,\widehat P_j)
    \label{eq:random-w2-lower-transfer}\\
    W_{2,\Gamma}^{2}(\widehat P_i,\widehat P_j)
    &\le L^{2}G_{ij}^{2}+2\tau_{ij}D_{ij}+\tau_{ij}^{2}.
    \label{eq:random-w2-upper-transfer}
  \end{align}
\end{proposition}

\begin{proof}
  Let $M=(m_{ab})$ be a transport matrix. Because the indexed characteristic
  and exposure clouds use the same atom masses, every nonnegative matrix with
  row sums $p_{ia}$ and column sums $p_{jb}$ is feasible for both the
  characteristic and exposure problems. Thus both Wasserstein costs minimize
  over the same coupling-matrix polytope.
  Write $Y_{ia}:=t_\Gamma(x_{ia})$ and
  $Z_{ia}:=\Gamma^{1/2}T_i(x_{ia},u_{ia})$, and define $Y_{jb},Z_{jb}$
  analogously. The indexed slack bounds imply
  $\norm{Z_{ia}-Y_{ia}}\le\tau_i$ and
  $\norm{Z_{jb}-Y_{jb}}\le\tau_j$.
  For each entry $(a,b)$, the bi-Lipschitz inequality gives
  \[
    L^{-2}d_\Omega(x_{ia},x_{jb})^2
    \le \norm{t_\Gamma(x_{ia})-t_\Gamma(x_{jb})}^2
    \le L^2d_\Omega(x_{ia},x_{jb})^2.
  \]
  The indexed slack vectors have norms at most $\tau_i$ and $\tau_j$, so the
  triangle inequality gives, entry by entry,
  \[
    \left|
    \norm{Z_{ia}-Z_{jb}}^2
    -\norm{t_\Gamma(x_{ia})-t_\Gamma(x_{jb})}^2
    \right|
    \le 2\tau_{ij}D_{ij}+\tau_{ij}^2.
  \]
  Combining the two entrywise displays, average against an arbitrary feasible
  $M$. Because every admissible $M$ is available in both problems, minimizing
  the lower and upper comparisons over the shared polytope gives
  equations~\ref{eq:random-w2-lower-transfer}--\ref{eq:random-w2-upper-transfer}.
\end{proof}

The finite-support proposition shows why matching atom masses matter: the
same transport weights can be evaluated before and after transmission, so
bounded distortion and slack carry characteristic-cloud distance into an
exposure-distance interval.
Economically, the interval is conditional on the declared values of $L$,
$\tau_{ij}$, and $D_{ij}$; permitting more transmission distortion or a
larger support bound weakens what a finite characteristic cloud can imply
about its indexed exposure cloud.
